\chapter{Introduction}

\begin{quote}
\textit{« Celui qui connaît les autres est instruit. Celui qui se connaît est sage. »} — Lao-Tseu
\end{quote}

Cette citation de Lao-Tseu résonne particulièrement avec ma démarche de Validation des Acquis de l'Expérience. Au-delà de l'instruction technique que j'ai acquise au fil des années, cette VAE représente un processus de connaissance approfondie de mes compétences et de leur transférabilité vers l'enseignement supérieur, dans le respect du cadre défini par l'INSPÉ \parencite{vademecum_vae_inspe}.

\section{Situation actuelle et motivations}

Cette année, je suis pleinement engagé en tant qu'enseignant en IUT, au sein du département Techniques de Commercialisation de l'IUT du Havre, où j'assure les enseignements de Ressources et Culture Numérique. Cette transition, effective depuis septembre 2025, marque une étape significative dans mon parcours professionnel, après quatre années d'enseignement des Sciences Numériques et Technologie (SNT) et de la spécialité Numérique et Sciences Informatiques (NSI) en lycée conformément aux programmes officiels \parencite{eduscol_nsi,nsi_programmes}.

Mon expérience s'est également enrichie par des interventions en vacations dans l'enseignement supérieur, notamment en troisième année de BUT, ce qui m'a permis d'appréhender les spécificités de l'ingénierie de formation pour des publics diversifiés. Cette évolution vers l'enseignement supérieur et la formation d'adultes constitue une étape significative de mon parcours professionnel et justifie pleinement ma démarche VAE pour le Master MEEF Pratiques et Ingénierie de la Formation (PIF) \parencite{unicaen}.

Mes motivations pour engager cette validation s'articulent autour de trois axes principaux. Premièrement, je recherche une légitimité académique en ingénierie de formation pour structurer mes pratiques. Deuxièmement, je prépare un projet de thèse que je souhaite débuter dans l'année à venir, centré sur l'impact des outils d'intelligence artificielle sur la conception de dispositifs de formation innovants et inclusifs \parencite{intelligence_artificielle_education}. Enfin, cette VAE s'inscrit dans une logique de professionnalisation en tant qu'expert en ingénierie pédagogique.

\section{Vision de l'ingénierie de formation à l'ère numérique}

Le domaine de la formation traverse aujourd'hui une période de mutations profondes. L'ingénierie pédagogique ne peut plus se concevoir sans intégrer pleinement les outils numériques et les nouvelles modalités d'apprentissage (hybride, distanciel, micro-learning). Je constate quotidiennement la nécessité de former non plus seulement des techniciens, mais des apprenants agiles, capables d'évoluer dans un environnement technologique en perpétuelle mutation \parencite{competences_21_siecle}.

Mon objectif en tant qu'expert en ingénierie de formation consiste à concevoir des dispositifs qui permettent aux apprenants de maîtriser les outils numériques pour en faire des leviers de leur développement professionnel. Cette approche holistique de la formation est d'autant plus cruciale avec l'avènement de l'IA générative. Ma participation à des groupes de travail nationaux sur l'IA en éducation (DNE) nourrit ma réflexion sur l'évolution nécessaire des modèles de formation.

Les enjeux futurs de notre discipline me paraissent multiples. L'intelligence artificielle occupe désormais une place centrale dans nos préoccupations pédagogiques, nécessitant une approche éthique et responsable de ces outils. Parallèlement, la sensibilisation à l'écologie numérique devient incontournable : nous devons former nos étudiants à une utilisation responsable des ressources numériques. Enfin, l'émergence de nouvelles méthodes d'évaluation permises par les outils numériques ouvre des perspectives pédagogiques prometteuses, particulièrement pour l'accompagnement des élèves à besoins particuliers \parencite{accessibilite_numerique}.

\section{Fil conducteur du parcours}

Mon évolution professionnelle suit une logique de progression cohérente : de l'expertise technique industrielle vers l'enseignement, puis vers l'ingénierie de formation en IUT, et enfin vers la recherche. Cette trajectoire illustre un passage progressif de la transmission de savoirs techniques à la conception de dispositifs de formation complexes \parencite{recherche_action_education}.

Ce qui fait sens dans mon parcours, c'est cette volonté permanente d'adapter l'ingénierie pédagogique aux évolutions technologiques et aux besoins spécifiques des apprenants. Mon engagement dans la veille technologique, l'expérimentation d'outils d'IA pour la formation et la prise en compte des enjeux éthiques témoignent d'une posture de praticien-chercheur en ingénierie de formation \parencite{perrenoud_pratique_reflexive}.

\section{Organisation de ce dossier}

Ce dossier s'organise en trois parties distinctes et complémentaires, conformément aux recommandations méthodologiques de l'INSPÉ \parencite{vademecum_vae_inspe} :

\subsection*{Partie 1 : Description du parcours}

Je présenterai mon parcours de candidat en détaillant :
\begin{itemize}[leftmargin=*]
    \item Ma formation initiale et continue
    \item Mes expériences professionnelles (industrie et enseignement)
    \item L'évolution de mes compétences techniques et pédagogiques
\end{itemize}

Cette première partie permettra d'établir le socle de connaissances et d'expériences sur lequel repose ma candidature.

\subsection*{Partie 2 : Description et analyse des expériences professionnelles}

J'analyserai mes expériences professionnelles les plus significatives, en mettant l'accent sur les situations qui ont contribué au développement de compétences en ingénierie de formation. Cette analyse démontrera la pertinence de mon parcours au regard des exigences du Master MEEF Pratiques et Ingénierie de la Formation (PIF) \parencite{rncp38152}.

\subsection*{Partie 3 : Correspondance avec le référentiel RNCP}

J'établirai les correspondances explicites entre mes acquis professionnels et le référentiel RNCP38152 \parencite{rncp38152}, démontrant que mon profil d'expert technique et pédagogique constitue un atout majeur pour l'ingénierie de formation.

\vspace{1cm}

Cette démarche VAE représente bien plus qu'une simple validation de compétences : elle constitue une étape fondamentale dans ma construction professionnelle, me permettant de donner du sens à mon parcours tout en préparant les défis pédagogiques de demain \parencite{schon_reflective_practitioner}.